\documentclass[10pt]{article}

% Rough manuscript draft. Replace the generic format and bibliography with the
% target venue's style once the submission venue is known.
\usepackage[margin=1in]{geometry}
\usepackage{booktabs}
\usepackage{tabularx}
\usepackage{amsmath}
\usepackage{microtype}
\usepackage{hyperref}
\usepackage{xcolor}

\newcommand{\todo}[1]{\textcolor{red}{[TODO: #1]}}
\newcommand{\taska}{Task A}
\newcommand{\taskb}{Task B}

\title{Parameter-Efficient Adaptation for Kannada--English Hate Speech Classification}
\author{\todo{Author names and affiliations}}
\date{}

\begin{document}
\maketitle

\begin{abstract}
Code-mixed social media text presents a difficult classification setting when
the dominant language is written in a non-standard Romanized form. We study
binary hate-speech detection and six-way hate-category classification for
Kannada--English YouTube comments, a setting in which Kannada is commonly
transliterated into Latin script (Kanglish). Starting from character and word
n-gram baselines and fine-tuned MuRIL encoders, we develop a parameter-efficient
decoder-language-model classifier. The classifier loads Gemma-4-12B in 4-bit
NF4 quantization, adapts selected attention and MLP projections with LoRA, and
predicts labels from a learned linear head over the final prompt-token
representation. On fixed stratified holdouts assembled from the released
training and labelled validation data, the recorded Gemma submissions achieve
macro-F1 scores of 0.8563 on the binary task and 0.6872 on the six-way task.
Subsequent holdout experiments report higher scores after task-adaptive
next-token pretraining (TAPT) and averaging independently adapted models, up to
0.8721 and 0.7056, respectively. We describe the training and evaluation
pipeline, its memory-aware implementation, and the distinction between
submitted predictions and later experiments. These results are promising
evidence for parameter-efficient domain adaptation on small code-mixed datasets;
they do not establish test-set performance or a general systems advantage.
\end{abstract}

\section{Introduction}
Online hate-speech moderation must handle language mixing, spelling variation,
and the limited labelled data available for many languages. These challenges
are pronounced for Kannada comments written in Latin characters and mixed with
English. A word-level representation may treat closely related spellings as
unrelated, while a pretrained model's tokenizer may fragment the transliterated
Kannada text.

This paper studies two HASTIKA-ICON-2026 classification tasks on Kannada--English
YouTube comments. \taska{} predicts \texttt{Hate} or \texttt{Non-Hate}; \taskb{}
assigns one of six categories: Gender, Geo-political, Others, Political,
Religion, or Violence. We follow a progression from sparse character and word
features, through MuRIL fine-tuning, to a quantized decoder language model with
parameter-efficient adaptation. We also study task-adaptive pretraining on
unlabelled in-domain text and probability averaging across independent runs.

The project began as a competition system and accumulated multiple generations
of experiments. We therefore distinguish the predictions retained as the
competition submissions from later research candidates. The submitted Gemma
systems were trained on all available labelled rows after model choices were
made on fixed holdouts. Later TAPT results are holdout measurements only; the
repository does not contain a scored test evaluation for them.

\paragraph{Contributions.} This draft makes three empirical contributions:
(1) it documents a controlled comparison of sparse, encoder, and parameter-
efficient decoder approaches for Romanized Kannada hate-speech classification;
(2) it describes a memory-conscious QLoRA and TAPT pipeline usable with limited
GPU memory; and (3) it records the effects and uncertainty of domain adaptation
and multi-run probability averaging on two related tasks. The implementation
details are engineering choices for this application rather than a claim of a
new general-purpose systems architecture.

\section{Task and Data}
The corpus consists of short Kannada--English YouTube comments. Kannada is
predominantly written in Roman characters in the data, and comments may include
English, emoji, URLs, mentions, repeated characters, and spelling variation.
The dataset paper and annotation protocol should be cited here
\cite{hastika_dataset}.

The labelled data used for the final fitting runs combine the original training
partition with a subsequently released labelled validation partition. After
the project's duplicate handling, the final fitting counts recorded in the
methods documentation are 7,193 rows for \taska{} and 3,532 for \taskb{}. The
fixed model-selection holdouts contain 1,079 and 530 rows, respectively. The
test inputs contain 806 binary-task comments and 395 six-way-task comments in
the released files. Test labels were unavailable to this project.

\begin{table}[t]
\centering
\small
\begin{tabularx}{\linewidth}{lXr}
\toprule
Task & Labels & Labelled rows for final fitting \\
\midrule
\taska{} & Hate, Non-Hate & 7,193 \\
\taskb{} & Gender, Geo-political, Others, Political, Religion, Violence & 3,532 \\
\bottomrule
\end{tabularx}
\caption{Task definitions and labelled data counts recorded for the final fits. Counts
reflect the project's preprocessing and duplicate handling; confirm them against the
official dataset release before submission.}
\label{tab:data}
\end{table}

\section{System Overview}
The codebase separates shared preprocessing and validation utilities from model
implementations, task-specific training code, experiment orchestration,
notebooks, and retained results. The supported pipeline has four stages:
input repair and normalization; optional unlabelled-text adaptation; supervised
training; and probability-based prediction packaging. Training and inference
preserve comment IDs, and the packaging utility checks that each required ID is
present exactly once and that all labels belong to the task vocabulary.

\subsection{Preprocessing}
The released CSVs contain text affected by UTF-8/Latin-1 mojibake. The shared
cleaner attempts to reverse this encoding issue, applies \texttt{ftfy}, decodes
HTML entities, replaces URLs and user mentions with normalized markers, and
collapses whitespace. For MuRIL, emoji can optionally be rendered as English
descriptions because of limited tokenizer coverage. Gemma experiments use the
native tokenizer and do not require this step.

For labelled training data, identical cleaned comments with agreeing labels are
collapsed to one row. Groups with conflicting labels are removed as ambiguous.
This reduces repeated weight in the loss and prevents identical text from
crossing a fold boundary. Prediction inputs retain every original ID, including
duplicate comments, to satisfy the data interface.

\subsection{Baselines and MuRIL}
The sparse baseline combines word and character TF--IDF n-grams with a
LinearSVC. Character features provide a strong reference for spelling and
morphological variation. The neural encoder line uses MuRIL with a pooled
classification head and, depending on the experiment, layer-wise learning-rate
decay, top-layer reinitialization, fast-gradient-method perturbations, and
exponential moving averages. For \taskb{}, class-weighted loss addresses the
imbalance between the six labels. Task-adaptive masked-language-model training
on Kannada text improved the earlier MuRIL system and motivated analogous
adaptation for the decoder model.

\subsection{Quantized decoder classifier}
The later classifier uses Gemma-4-12B as a representation encoder rather than
generating label strings. The backbone is loaded in 4-bit NF4 with double
quantization and fp16 computation. LoRA adapters (rank 16, alpha 32, dropout
0.05) are applied to the attention projections
\texttt{q\_proj}, \texttt{k\_proj}, \texttt{v\_proj}, \texttt{o\_proj} and MLP
projections \texttt{gate\_proj}, \texttt{up\_proj}, and \texttt{down\_proj}.
Gradient checkpointing reduces activation storage. The large token-generation
head is omitted during classification.

Each example is placed in a short task-specific prompt ending in an answer cue
(for example, \texttt{Category:}). Let $h_i$ be the final hidden state at the
last non-padding prompt token for example $i$. The classifier computes
\[
  z_i = W\,\mathrm{Dropout}(h_i) + b,
  \qquad p_i = \mathrm{softmax}(z_i),
\]
where $W$ and $b$ are learned from scratch. Training uses weighted cross-entropy
with label smoothing 0.05. The standard LoRA and head learning rates are
$10^{-4}$ and $5\times10^{-4}$, respectively; optimization uses AdamW, a cosine
schedule with warmup, and an effective batch size of 16 through gradient
accumulation. The task-specific epoch counts are three for \taska{} and four
for \taskb{} in the recorded submission recipes.

\subsection{Task-adaptive pretraining}
The decoder TAPT stage continues next-token training on unlabelled in-domain
comments and permitted external Kannada text, using LoRA on the same decoder
projections as the classifier. The resulting adapter initializes the supervised
stage. The code deduplicates text and can restrict TAPT inputs through explicit
corpus paths; evaluation/test text must not be included in a clean inductive
comparison. For Gemma, the large vocabulary projection is applied to chunks of
hidden states so that the full token-logit tensor need not be held in memory at
once. This implementation detail was added after an earlier TAPT attempt
exceeded available GPU memory.

\subsection{Training safeguards and ensembling}
Variable-length examples are grouped into length-similar minibatches to reduce
padding. Mixed precision, gradient accumulation, gradient clipping, and
non-finite-loss checks support training on a single accelerator. A collapse
guard checks the first-epoch predictions and terminates a run if it predicts one
class for nearly all sampled training examples despite a non-degenerate class
prior. Failed or collapsed runs are excluded from probability ensembles.

For ensembles, each member writes probabilities in a documented class order.
Members are averaged before the final argmax. The final combination utility
weights each saved probability matrix by its number of constituent models,
checks dimensions and probability sums, and rejects collapsed members. Earlier
MuRIL/SVM ensembles used aligned out-of-fold predictions to estimate blend
weights; the Gemma recipes use equal weighting rather than holdout-tuned weights.

\section{Experimental Protocol}
The primary metric is macro-averaged F1, with accuracy and per-class F1 as
secondary measures. Macro-F1 gives each category equal weight, which is
important for \taskb{}'s minority labels. The fixed holdouts were generated
with stratification and a shared split seed, and their IDs are fingerprinted in
the repository so notebook runs evaluate the same examples. Older MuRIL/SVM
experiments also use five-fold out-of-fold predictions. OOF, fixed-holdout,
and competition scores are reported separately because they are not directly
comparable.

For paired comparisons, the project resamples holdout examples with replacement
and reports a percentile interval for the macro-F1 difference. The holdouts are
small, especially for the minority categories, so intervals can be wide and
small score differences are uncertain. Multiple model seeds are used to expose
optimization variance; seed selection based on test predictions is not used.

The final submitted Gemma runs fit on the combined labelled data after choices
were made on the fixed holdouts. Later TAPT runs used matching holdout designs
to compare adapted models, but their results must not be presented as test-set
scores. The experiment history includes some exploratory transductive
experiments; those are excluded from the clean method described here and should
be individually identified if discussed.

\section{Results}
Table~\ref{tab:results} summarizes representative measurements documented in
the repository. Within each task, the principal comparison uses the same fixed
holdout. These figures are draft-level until each score is cross-checked against
its notebook output and run log.

\begin{table}[t]
\centering
\small
\begin{tabular}{llrl}
\toprule
Task & System & Holdout macro-F1 & Status \\
\midrule
\taska{} & Character SVM & 0.8117 & Baseline \\
\taska{} & Gemma QLoRA, one seed & 0.8563 & Recorded submission recipe \\
\taska{} & Gemma TAPT ensemble, later run & 0.8702 & Post-submission holdout \\
\taska{} & TAPT and R-Drop mixture, later run & 0.8721 & Post-submission holdout \\
\midrule
\taskb{} & MuRIL & 0.6151 & Earlier holdout reference \\
\taskb{} & Gemma QLoRA, four-model recipe & 0.6872 & Recorded submission recipe \\
\taskb{} & Independent TAPT pair, later run & 0.7056 & Post-submission holdout \\
\bottomrule
\end{tabular}
\caption{Representative macro-F1 measurements. Rows marked as submission recipes
describe the model-selection evidence for the systems whose prediction artifacts were
submitted; later TAPT rows are follow-up holdout experiments, not scored test results.
Baseline protocols should be verified for exact split matching before publication.}
\label{tab:results}
\end{table}

The recorded Gemma submission recipes improve substantially over the earlier
MuRIL/SVM era in the fixed-holdout comparisons. Later experiments suggest that
TAPT can provide additional gains, but the results also reveal sensitivity to
the initialization and seed: two Task B TAPT pairs in the ledger differ by
several macro-F1 points. This variation argues for reporting ensembles and
uncertainty rather than relying on a single best run. Because these later scores
were used during iterative experimentation on fixed holdouts, they should be
treated as development evidence and confirmed on an untouched test set or a
new independent split before making strong generalization claims.

\section{Discussion}
Three observations motivate the system design. First, character n-gram models
remain useful baselines for transliterated text, where surface variation can
fragment word-level evidence. Second, broad pretrained decoder models can
produce stronger representations than the earlier encoder and sparse models in
this setting, even when only a small LoRA parameter set and a linear head are
trained. Third, domain adaptation and averaging appear complementary: TAPT
teaches the model the corpus's spelling and code-mixing patterns, while
independent fits reduce sensitivity to a particular initialization.

These observations are task- and corpus-specific. They do not establish that a
larger decoder is always preferable, nor do they isolate model scale from
pretraining data, tokenizer behavior, and optimization. A stronger systems
study would measure throughput, peak memory, energy or cost, and inference
latency under controlled hardware and matched quality targets. The current
repository records some runtime and memory diagnostics, but it was not designed
as a controlled hardware benchmark.

\section{Limitations and Ethics}
The released test labels are not available in this repository, so we cannot
report final test performance. The fixed holdouts are relatively small, and the
experiment history contains many adaptive comparisons; the best observed score
is likely optimistic. Several candidate models and experiment summaries are
stored as notebooks or logs rather than as complete, locally retained
checkpoints and prediction matrices. Exact reproducibility therefore requires
the original model access, compatible software versions, hardware, and the
relevant notebook outputs.

Hate-speech labels encode annotation decisions and may be ambiguous without
conversation context. The project drops exact duplicate texts with conflicting
labels for training, but this does not resolve broader label subjectivity.
Automated hate-speech classification can cause harm through false positives,
false negatives, or uneven performance across communities and topics. The
system should be treated as a research artifact and not as a stand-alone
moderation decision-maker. A responsible deployment would require human review,
community-informed error analysis, and evaluation across dialects, demographic
groups, and changing language use.

\section{Conclusion}
We described a small-data classification pipeline for Kannada--English
Romanized social-media text, progressing from character n-gram and MuRIL
baselines to 4-bit QLoRA classification with Gemma-4-12B. The retained
competition submissions provide a reproducible historical point of comparison,
while later holdout experiments indicate that LoRA-based TAPT and independent
model averaging may improve both tasks. The present evidence supports further
evaluation of parameter-efficient domain adaptation; it does not yet establish
general systems efficiency or performance on the unlabelled test set.

\section*{Reproducibility Checklist}
Before submission, complete the following items:
\begin{itemize}
  \item Confirm dataset and split counts, labels, and duplicate policy against the release.
  \item Identify exact commits, notebook versions, model revisions, and software versions.
  \item Reconcile README, final-submission, methods, and experiment-ledger status.
  \item Verify every reported metric, interval, and baseline split from saved outputs.
  \item Report GPU model, per-run wall time, peak memory, and inference cost if making systems claims.
  \item Add a complete related-work discussion and verified bibliography.
  \item Confirm licenses, external-corpus provenance, and model access conditions.
\end{itemize}

\bibliographystyle{plain}
\begin{thebibliography}{9}
\bibitem{hastika_dataset}
\todo{Dataset paper authors}.
\newblock \todo{Exact title of the HASTIKA dataset paper}.
\newblock \todo{Venue and year}. Replace with the verified citation from
\texttt{references/hastika\_dataset\_paper.pdf}.

\bibitem{muril}
\todo{MuRIL paper authors}.
\newblock \todo{Exact MuRIL paper title, venue, and year}.

\bibitem{qlora}
\todo{QLoRA paper authors}.
\newblock \todo{Exact QLoRA paper title, venue, and year}.

\bibitem{lora}
\todo{LoRA paper authors}.
\newblock \todo{Exact LoRA paper title, venue, and year}.
\end{thebibliography}

\end{document}
